\documentclass[10pt,letterpaper]{article}

\usepackage[letterpaper,margin=0.70in,headheight=22pt,headsep=16pt,footskip=24pt]{geometry}
\usepackage[T1]{fontenc}
\usepackage[utf8]{inputenc}
\usepackage{lmodern}
\usepackage{microtype}
\usepackage{xcolor}
\usepackage{graphicx}
\usepackage{booktabs}
\usepackage{array}
\usepackage{tabularx}
\usepackage{longtable}
\usepackage{amsmath,amssymb}
\usepackage{tikz}
\usetikzlibrary{arrows.meta,positioning,calc,shapes.geometric,decorations.pathreplacing,fit}
\usepackage{enumitem}
\usepackage{fancyhdr}
\usepackage{titlesec}
\usepackage{caption}
\usepackage{listings}
\usepackage{xurl}
\usepackage[hidelinks]{hyperref}
\usepackage{lastpage}
\usepackage{float}
\usepackage{placeins}
\usepackage{needspace}

\definecolor{TrawlAmber}{HTML}{A66A00}
\definecolor{TrawlGray}{HTML}{555555}
\definecolor{TrawlBlue}{HTML}{144E8C}
\definecolor{LightGray}{HTML}{F3F3F3}
\definecolor{MidGray}{HTML}{D9D9D9}
\definecolor{DarkText}{HTML}{1A1A1A}
\definecolor{GoodGreen}{HTML}{2B7A2B}
\definecolor{WarnOrange}{HTML}{A65E00}

\renewcommand{\familydefault}{\sfdefault}
\setlength{\parindent}{0pt}
\setlength{\parskip}{5pt}
\renewcommand{\arraystretch}{1.15}
\sloppy
\setlist[itemize]{leftmargin=1.2em,itemsep=2pt,topsep=2pt}
\setlist[enumerate]{leftmargin=1.4em,itemsep=2pt,topsep=2pt}

\titleformat{\section}{\Large\bfseries\color{DarkText}}{}{0pt}{}[\vspace{-3pt}\color{TrawlAmber}\titlerule]
\titleformat{\subsection}{\large\bfseries\color{DarkText}}{}{0pt}{}
\titleformat{\subsubsection}{\normalsize\bfseries\color{DarkText}}{}{0pt}{}

\newcommand{\code}[1]{\texttt{\detokenize{#1}}}
\newcolumntype{Y}{>{\raggedright\arraybackslash}X}
\newcolumntype{L}[1]{>{\raggedright\arraybackslash}p{#1}}
\captionsetup{font=small,labelfont=bf}

\lstdefinestyle{shell}{
  basicstyle=\ttfamily\footnotesize,
  breaklines=true,
  columns=fullflexible,
  frame=single,
  rulecolor=\color{MidGray},
  backgroundcolor=\color{LightGray},
  keepspaces=true,
  showstringspaces=false
}
\lstset{style=shell}

\newcommand{\notebox}[1]{%
  \vspace{4pt}\noindent\fcolorbox{TrawlBlue}{LightGray}{%
  \begin{minipage}{0.965\linewidth}%
  \textbf{Note:} #1%
  \end{minipage}}\vspace{4pt}}

\newcommand{\importantbox}[1]{%
  \vspace{4pt}\noindent\fcolorbox{TrawlAmber}{LightGray}{%
  \begin{minipage}{0.965\linewidth}%
  \textbf{Important:} #1%
  \end{minipage}}\vspace{4pt}}

\newcommand{\field}[1]{\texttt{\detokenize{#1}}}

\pagestyle{fancy}
\fancyhf{}
\fancyhead[L]{\small Datatrawl-UG001}
\fancyhead[C]{\small Storage-Safe Archive Trawling Engine v1.0.0}
\fancyhead[R]{\small User Guide}
\fancyfoot[L]{\small Datatrawl-UG001 v1.3}
\fancyfoot[C]{\small West Virginia University}
\fancyfoot[R]{\small Page \thepage\ of \pageref{LastPage}}
\renewcommand{\headrulewidth}{0.8pt}
\renewcommand{\headrule}{\hbox to\headwidth{\color{TrawlAmber}\leaders\hrule height \headrulewidth\hfill}}
\renewcommand{\footrulewidth}{0.8pt}
\renewcommand{\footrule}{\hbox to\headwidth{\color{TrawlAmber}\leaders\hrule height \footrulewidth\hfill}}

\title{datatrawl Storage-Safe Archive Trawling Engine v1.0.0\\User Guide}
\author{West Virginia University}
\date{Datatrawl-UG001 v1.3 \\ July 10, 2026}
\hypersetup{
  pdftitle={datatrawl Storage-Safe Archive Trawling Engine v1.0.0 -- User Guide},
  pdfauthor={West Virginia University Radio Astronomy Instrumentation Lab (RAIL)},
  pdfsubject={Datatrawl-UG001 v1.3},
  pdfkeywords={datatrawl, radio astronomy, archive analysis, user guide}
}

\begin{document}
\thispagestyle{fancy}
{\sffamily
\begin{minipage}[t]{0.39\textwidth}
  {\fontsize{24}{26}\selectfont\bfseries\color{TrawlAmber}datatrawl}\par
  {\small Datatrawl-UG001 v1.3, July 10, 2026}
\end{minipage}%
\begin{minipage}[t]{0.61\textwidth}
  \raggedleft
  {\large\bfseries User Guide}\par
  {\small Install, plan, run, resume.}\par
  {\small West Virginia University Radio Astronomy Instrumentation Lab (RAIL)}
\end{minipage}
}
\vspace{6pt}
{\color{TrawlAmber}\hrule height 1.2pt}
\vspace{8pt}

\section{The Model in One Paragraph}
Long archive scans are limited by scratch space and interrupted computing
sessions. We address these constraints by processing one primary file as one
unit and checkpointing a small analysis product. A run combines four choices:
a \textbf{telescope} supplies geometry and metadata, a \textbf{source} lists
and stages files, a \textbf{reader} converts one staged file into arrays, and
an \textbf{analyzer} performs the science calculation. The engine fetches,
analyzes, and removes each unit while keeping the number of staged files within
the configured bound. Use \code{datatrawl doctor} to check the four choices
before starting a long scan.

\section{Installation}
\begin{lstlisting}
git clone https://github.com/WVURAIL/datatrawl.git
cd datatrawl
python -m venv .venv && source .venv/bin/activate
pip install -e ".[dev,examples]"   # offline/local development
# For archive/CANFAR use instead:
#   pip install -e ".[survey,examples,dev]"
#   cadc-get-cert -u <your_cadc_username>
#   datatrail config init --site canfar
datatrawl --version
pytest -q
\end{lstlisting}
The \code{[cadc]} extra supports CADC fetching from an existing inventory. The
\code{[survey]} extra adds the CADC clients and
\code{datatrail-cli>=0.11.0} for inventory construction. The \code{[gpu]}
extra does not pin CuPy because the installed CuPy build must match the session
image (see \S\ref{sec:cupy}).

\section{Discovering What Is Available}
\begin{lstlisting}
datatrawl list                 # everything at a glance
datatrawl list telescopes      # ready vs. stub
datatrawl list readers
datatrawl list analyzers
datatrawl list sources
\end{lstlisting}
Next, check the combination that will be used for the scan:
\begin{lstlisting}
datatrawl doctor --telescope chime --source cadc-datatrail \
                 --reader chime-baseband --analyzer spectrum
\end{lstlisting}
\code{doctor} reports the installed components, configuration, credentials,
and readiness of the selected combination.

When the archive layout is not yet known, first build a discovery map. The
\code{survey --scopes-only} mode lists datasets across scopes without
enumerating events or downloading files. It writes the result to
\code{scopes.jsonl}, or to
\code{scopes-<name>.jsonl} with \code{--name}, so successive recons do not
overwrite one another:
\begin{lstlisting}
datatrawl survey --telescope gbo --scopes-only              # this telescope's scopes
datatrawl survey --scopes-only                              # every scope datatrail sees
datatrawl survey --telescope gbo --scopes-only \
    --scope gbo.acquisition.processed --match gain --expand --name gains
\end{lstlisting}
\Needspace{11\baselineskip}
\code{--telescope} restricts discovery to live scopes whose first name
component matches the telescope. These scopes come from Datatrail rather than
the instrument YAML, because the YAML only declares scopes used by the event
survey. Omit \code{--telescope} to inspect every visible scope.
\code{--match} keeps names containing all comma-separated terms.
\code{--expand} opens each retained dataset by one level, so a container such
as \code{complex_gains} contributes its children as rows with a
\field{parent} column. Each row can then be inspected with
\code{datatrail ps <scope> <dataset> -s}. If Datatrail does not answer for a
scope, the command identifies that scope and marks the map INCOMPLETE.
\notebox{Discovery maps names and containers. The event survey enumerates only
event-keyed products. For a timestamped product such as a calibration
acquisition, use \code{--scopes-only --expand}, then list the selected dataset
with \code{DATATRAIL.files()}. See
\code{docs/DATATRAIL_BOUNDARY.md} for the division of responsibility.}

\section{Building an Inventory}
\code{survey} converts the archive listing into the units required by one
analysis. It writes a named inventory directory, using
\code{data/<telescope>-fid<freq_ids>/} by default or the value passed to
\code{--name}. The inventory metadata records the telescope, source, and
reader, so a later scan only needs the inventory and analyzer:
\begin{lstlisting}
datatrawl survey \
  --telescope chime --source cadc-datatrail \
  --scope chime.event.baseband.raw \
  --freq-ids 844 --max-events 5 --name chime-spectrum-844

datatrawl explore --name chime-spectrum-844
\end{lstlisting}
The \code{--max-events 5} setting bounds this first survey. Survey verifies
candidate files but does not bulk-download them. The inventory remains on disk
between commands. During a scan, the current engine loads the rows selected for
one product into memory before staging begins. Memory use therefore grows with
the selected inventory size.

The reader defines the \emph{archive file shape}: the files expected for one
event and their names (\field{Reader.survey_files}). Survey uses the
telescope's canonical reader unless \code{--reader} overrides it. Each
verified row records its \field{name}; a later scan stages that recorded file
instead of reconstructing the naming rule. To survey a non-baseband product,
load a reader that defines its archive shape with \code{--plugin}; see
\code{docs/ADDING_A_READER.md}.

\section{Running a Scan}
\begin{lstlisting}
datatrawl scan \
  --name chime-spectrum-844 --analyzer spectrum --select 844 \
  --max-frames-per-file 5
\end{lstlisting}
This command limits each file to five frames and is intended as a smoke test.
For an uncapped run, use a new \code{--out} or remove the smoke-test product.
The spectrum analyzer rejects a resume when the cap differs.
\begin{itemize}
  \item \code{--select} is interpreted by the analyzer. For the spectrum
        analyzer: explicit \field{freq_id}s such as \code{844},
        \code{614,706}, or ranges like \code{506-552}. Event-oriented
        analyzers take \code{events:349382977[,...]}, and a
        \field{plan_runs} may return the dict form
        \code{{"events": [...], "freq_ids": ...}}; filters are exact and
        ANDed. A malformed selection reports the accepted grammar. The
        analyzer's \field{plan_runs} decides product fan-out;
        the built-in spectrum analyzer creates one independent product per
        selected \field{freq_id}.
  \item \code{--checkpoint-every N} (default 50) sets the checkpoint
        interval. Files consumed after the most recent completed checkpoint
        may need to be processed again after an interruption.
  \item \code{--download-workers} sets the number of fetch threads; concurrent
        fetches also require more than one staging slot.
        \code{--max-staged-files} caps fetched or in-progress files on scratch
        and permits fetch/analyze overlap. Analyzer calls remain serial on the
        engine's main thread. The engine refuses either value above 1 when the
        analyzer declares \field{requires_in_order}.
  \item \code{--plugin yourpkg.module} loads an external plugin explicitly;
        installed packages advertising the \code{datatrawl.plugins}
        entry-point group are found automatically.
\end{itemize}

\importantbox{At the defaults, one staged or in-progress file occupies scratch.
If \code{--max-staged-files} is raised, the configured number is the file-count
bound. Without \code{--tmp-dir}, each invocation receives a unique staging
directory. An explicit \code{--tmp-dir} is used as-is, so concurrent scans
require distinct paths. The engine removes the files it stages; use a scratch
directory that does not contain source data or other required files.}

\section{Interrupt and Resume}
A session timeout, preemption, or \code{Ctrl-C} may stop a scan. Repeat the
same command with the same output to continue. The analyzer reloads the last
completed checkpoint and reports its committed unit keys; the engine skips
matching units. Files consumed after that checkpoint may be processed again.
The supplied analyzer base class and spectrum analyzer write a temporary
product in the output directory and replace the prior checkpoint only after the
write completes. An external analyzer owns the equivalent recovery behavior.
If a changed parameter would alter product meaning, a conforming analyzer
rejects the resume. Use a new output location when intentionally changing the
configuration.

\Needspace{7\baselineskip}
\section{CANFAR / CADC Notes}
On CANFAR, install \code{[survey]} when building an inventory; this extra
includes the CADC clients. Use \code{[cadc]} alone when scanning an existing
inventory without Datatrail discovery. Obtain a valid CADC proxy certificate,
then run \code{datatrawl doctor --source cadc-datatrail ...} to check the
credentials and Datatrail interface. Datatrail remains the archive authority
and datatrawl reads from it. The boundary is defined in
\code{docs/DATATRAIL_BOUNDARY.md}. For a long scan, choose
\code{--checkpoint-every} by balancing repeated work after interruption
against product-write overhead.

\section{GPU Sessions and CuPy}\label{sec:cupy}
GPU analyzers use CuPy. The CuPy build must match the CUDA runtime in the
session image, so datatrawl does not pin it as a package dependency. On a GPU
node:
\begin{lstlisting}
datatrawl setup-cupy            # report the image's CUDA + the matching wheel
datatrawl setup-cupy --install  # install cupy-cudaXXx into the active venv
\end{lstlisting}
Without \code{--gpu}, the built-in spectrum analyzer uses NumPy. With
\code{--gpu}, it requires CuPy; if CuPy cannot be imported, the scan reports
the setup command and stops. The CPU and GPU paths implement the same
calculation, although floating-point roundoff may differ.

\section{Writing an Analyzer}
The analyzer performs the analysis and owns the product contract. Implement
the lifecycle and register the class:
\begin{lstlisting}
class MyAnalyzer(Analyzer):
    info = PluginInfo(name="my-analyzer", ...)
    requires_in_order = False   # True for any order-dependent product

    def resume(self, path, ctx):      # load existing product; raise if
        ...                           # present-but-incompatible
    def processed_keys(self):         # unit keys already in the product
    def begin(self, ctx, first_meta): # once, at the first new file
    def consume_file(self, arrays, meta):  # accumulate; return n frames
    def save(self, path):             # persist product, keys, and provenance
    def summary(self):                # one status line
\end{lstlisting}
Two optional hooks define selection behavior. \field{resolve_selection}
interprets \code{--select}, and \field{plan_runs} divides that selection into
independent runs. Each run receives a new analyzer instance and product.
Analyzer calls occur on the engine's main thread. The complete lifecycle and
inventory-metadata contract are described in
\code{docs/ADDING_AN_ANALYZER.md} and the corresponding reader, source, and
telescope guides.

\Needspace{12\baselineskip}
The repository includes two worked patterns. For \emph{per-event fan-out},
\field{plan_runs} returns one
\code{{"events": [ev], "freq_ids": spec}} sub-selection per event, so each
event becomes its own resumable product
(\code{ev<event>[_<freq_ids>].npz}); both archive and local sources
understand this dictionary. For a \emph{companion side-load}, an analyzer that
needs calibration gains or flags loads a lookup once in \field{begin}, then
selects the companion from \field{meta["event"]} for each primary file. The
analyzer records the companion identity in its product and rejects a resume if
the lookup changes. Day-keyed companion archives can instead resolve with
\code{DATATRAIL.files(scope, day)}, imported from
\code{datatrawl.plugins.sources} --- the programmatic
\code{datatrail ps -s}. The runnable, CLI-tested reference is
\code{examples/per_event_companions.py}; the patterns are written out in
\code{docs/ADDING_AN_ANALYZER.md}.

\importantbox{If the product depends on consume order, including a running
statistic, trailing statistic, or CFAR baseline, set
\field{requires_in_order = True}. The engine then rejects either parallel or
prefetch setting above 1.}

\Needspace{0.20\textheight}
\section{Troubleshooting}
\begin{tabularx}{\linewidth}{L{0.34\linewidth} Y}
\toprule
\textbf{Symptom} & \textbf{First move} \\
\midrule
``How do I even start?'' & \code{datatrawl list}, then \code{datatrawl doctor}
with your four choices. \\
Scan refuses to start & Read the doctor-style preflight error: missing extra,
missing credential, or a reader/telescope mismatch. \\
Resume refused & The existing product was built with different parameters;
the analyzer prevents those configurations from being combined. Rebuild in a
new output location. \\
Parallel flags refused & The analyzer declares \field{requires_in_order};
run at the defaults. \\
Anything else & \code{docs/TROUBLESHOOTING.md} in the repository. \\
\bottomrule
\end{tabularx}

\section{Related Documents}
\textbf{Datatrawl-DS001} (product specification: contracts, limits,
interfaces); the \code{docs/ADDING_*.md} extension guides;
\code{docs/DATATRAIL_BOUNDARY.md}; and, for the RAIL detector built on this
engine, \textbf{PilotProxy-DS001/UG001}.

\Needspace{0.24\textheight}
\section{Revision History}
\begin{tabularx}{\linewidth}{L{0.12\linewidth} L{0.20\linewidth} Y}
\toprule
\textbf{Rev} & \textbf{Date} & \textbf{Changes} \\
\midrule
v1.0 & July 1, 2026 & Initial release, documenting engine v0.1.0. \\
v1.1 & July 3, 2026 & Engine v0.2.0: recon (\code{--scopes-only} with
\code{--match}/\code{--expand}/\code{--name}, telescope-scoped), event
selection grammar, reader-owned archive file shape (\code{survey --reader}),
\code{DATATRAIL.files()}, per-event fan-out and companion patterns. \\
v1.2 & July 3, 2026 & Version designation: engine v0.2.0 redesignated v1.0.0
under the 1.0 stability contract (public CLI, plugin contracts,
product/inventory formats); no functional changes. External analyzer example
moved to the \code{datatrawl-analyzer-template} repository. \\
v1.3 & July 10, 2026 & Documentation corrections: bounded prefetch and
scratch isolation, selection-sized inventory memory, analyzer-owned
provenance and resume requirements, GPU behavior, survey dependencies,
pagination, and PDF metadata; no engine behavior changes. \\
\bottomrule
\end{tabularx}

\end{document}
